\documentclass[../main_full_simulink.tex]{subfiles}

\begin{document}

\section{Simulink Reproduction of the Closed Loop}\label{sec:simulink}

The frequency-domain design of Tasks~1--3 is carried out entirely on transfer
functions and state-space objects in MATLAB. As an independent cross-check I
rebuilt the same closed loop as a Simulink block diagram,
\texttt{models/hm3\_closed\_loop.slx}, and verified that the two agree in the
time domain to solver tolerance. The MATLAB scripts remain the single source of
truth: every gain, filter and plant matrix is computed by
\texttt{main\_task1/2/3.m} and exported to the base workspace by
\texttt{init\_simulink\_hm3.m}; the model only \emph{references} those variables,
so nothing is retyped by hand and re-running the initialisation retunes the
model automatically.

\subsection{Model architecture}

The loop has four players: a \textbf{PD controller} (a single Gain block acting
on the measurement vector), an \textbf{actuator path} (the bending notch in
series with the second-order TVC servo and its Pad\'e delay --- or nothing, in
the rigid Task~1 case), a \textbf{plant} (one State-Space block, $4$-state rigid
or $6$-state flexible), and a \textbf{wind disturbance} $\alpha_w(t)$ supplied by
a \texttt{From Workspace} block. Tasks~1 and~2--3 differ \emph{only} in the
actuator path and the plant matrices, so both live inside a single
\textbf{Variant Subsystem} with two choices; the workspace variable
\texttt{task}~$\in\{1,2,3\}$ (pushed by \texttt{init\_simulink\_hm3}) selects the
active one at simulation time, while the controller, wind and logging are shared
and built once (Figure~\ref{fig:slx_root}).

\begin{figure}[H]
\centering
\includegraphics[width=0.92\textwidth]{hm3_closed_loop_root.png}
\caption{Root of the Simulink mirror: the PD controller (\texttt{Controller\_PD}),
the variant plant and actuator block (\texttt{Plant\_and\_Actuator}), the wind
input (\texttt{wind\_ts}) and the signal logging. The measurement bus feeds the
controller, whose output \texttt{u\_pd} drives the plant; the demultiplexed
states close the loop and are logged.}
\label{fig:slx_root}
\end{figure}

The measurement vector fed back to the controller is
$\mathbf{y}_\mathrm{meas} = [\theta_m,\ \dot\theta_m,\ z_m,\ \dot z_m]$, i.e.\ the
INS outputs, which in the flexible plant are contaminated by the bending
coordinate $\eta$ through the leakage gains $\phi_{1,ins},\sigma_{1,ins}$ of
Table~\ref{tab:params}. The plot vector $[\theta,\ z,\ \dot z]$ collects the true
states and is kept out of the loop, so the pitch history is displayed without the
bending contribution. The two interchangeable choices held in the variant
subsystem \texttt{Plant\_and\_Actuator} are shown in
Figure~\ref{fig:slx_choices}.

\begin{figure}[H]
\centering
\begin{subfigure}{0.9\textwidth}\centering
  \includegraphics[width=\linewidth]{hm3_closed_loop_rigid.png}
  \caption{\texttt{Rigid\_Ideal} (Task~1): rigid plant, ideal actuator ($\delta = u_{pd}$).}
  \label{fig:slx_rigid}
\end{subfigure}

\vspace{0.9em}
\begin{subfigure}{0.9\textwidth}\centering
  \includegraphics[width=\linewidth]{hm3_closed_loop_full.png}
  \caption{\texttt{Full\_TVC\_Notch} (Tasks~2--3): notch, TVC servo with Pad\'e delay, flexible plant.}
  \label{fig:slx_full}
\end{subfigure}
\caption{The two plant and actuator choices held in the variant subsystem; the
workspace variable \texttt{task} activates one at simulation time.}
\label{fig:slx_choices}
\end{figure}

\subsection{Validation against the scripts}

Each run drives the block diagram and the script loop with the \emph{same} wind
history: \texttt{run\_simulink\_closed\_loop.m} assembles the analytic
closed-loop transfer $T(s)$ with \texttt{assemble\_loop} and integrates it with
\texttt{simulate\_gust\_response}, then simulates \texttt{hm3\_closed\_loop.slx}
over the same horizon and overlays the two responses
(Figure~\ref{fig:slx_val}). The block diagram is integrated with \texttt{ode45}
while the script response is a transfer-function simulation, so the only expected
difference is the numerical integration error.

\begin{figure}[H]
\centering
\includegraphics[width=0.9\textwidth]{task1_simulink_vs_script.png}

\vspace{0.6em}
\includegraphics[width=0.9\textwidth]{task2_simulink_vs_script.png}
\caption{Simulink (dashed) versus script (solid) closed-loop response to the
sizing $1$-cosine gust: pitch $\theta$, lateral drift $z$ and TVC deflection
$\delta$. Top: Task~1 (rigid plant, ideal actuator); bottom: Task~2 (flexible
plant, TVC servo, Pad\'e delay and bending notch). The traces are
indistinguishable at plot scale.}
\label{fig:slx_val}
\end{figure}

The peak absolute discrepancies over the whole transient,
Table~\ref{tab:slx_agreement}, stay at the $10^{-6}$--$10^{-7}$~rad level in pitch
and TVC angle and below a millimetre in drift --- i.e.\ at solver tolerance ---
for the sizing gust, the stronger sustained-wind profile, and the worst-case
Task~3 corner (V3, where the frozen controller still keeps the perturbed plant
stable). The block diagram therefore reproduces the classical design point by
point, and the same \texttt{task} switch that selects the plant in the model
mirrors the progression of the three tasks in the report.

\begin{table}[H]
\centering
\caption{Peak absolute Simulink-vs-script discrepancy over the transient.}
\label{tab:slx_agreement}
\begin{tabular}{lccc}
\toprule
Case & $\max|\Delta\theta|$ [rad] & $\max|\Delta z|$ [m] & $\max|\Delta\delta|$ [rad] \\
\midrule
Task~1 --- gust        & $3.2\times10^{-6}$ & $8.8\times10^{-4}$ & $5.4\times10^{-6}$ \\
Task~2 --- gust        & $1.7\times10^{-7}$ & $1.0\times10^{-5}$ & $6.2\times10^{-7}$ \\
Task~2 --- strong wind & $8.7\times10^{-7}$ & $7.7\times10^{-5}$ & $1.8\times10^{-5}$ \\
Task~3 --- corner V3   & $3.7\times10^{-7}$ & $2.3\times10^{-5}$ & $9.4\times10^{-7}$ \\
\bottomrule
\end{tabular}
\end{table}

\subsection{Validation driver}

Listing~\ref{lst:hm3_simrun} is the core of the validation driver: it populates
the workspace, builds the analytic baseline with the \emph{same} controller and
plant the model embeds, simulates the block diagram and extracts the logged
signals for the overlay. The plotting boilerplate and the full block-by-block
construction of \texttt{hm3\_closed\_loop.slx} are omitted here and documented in
\texttt{models/SIMULINK\_GUIDE.md}.

\lstinputlisting[style=matlab, firstnumber=31, firstline=31, lastline=57,
  caption={Core of the Simulink validation driver (\texttt{run\_simulink\_closed\_loop.m}).},
  label=lst:hm3_simrun]{../run_simulink_closed_loop.m}

\end{document}
